\section{Analyse des résultats de BERTopic}

Le modèle final distingue 17 topics et attribue un topic à \(61\,419\) des \(61\,432\) segments du corpus ; seuls 13 segments conservent l'étiquette \texttt{-1}. Cette couverture presque complète rend possible une lecture d'ensemble, mais elle ne supprime pas les réserves formulées lors de la validation : \(46\,952\) affectations finales, soit 76,4\,\%, proviennent de la réattribution des documents initialement atypiques. Les effectifs décrivent donc principalement une allocation par proximité aux 17 noyaux sémantiques du modèle, et non la taille de 17 amas denses directement découverts par HDBSCAN.

Un topic BERTopic ne constitue pas un thème littéraire déjà nommé. Il rassemble des segments proches dans l'espace sémantique, puis les caractérise par un vocabulaire pondéré. Son interprétation repose ici sur trois indices complémentaires : les quinze termes les mieux classés par le c-TF-IDF, les documents représentatifs et la distribution des segments entre les romans. Les intitulés employés dans cette section sont par conséquent des étiquettes de lecture ; ils synthétisent le contenu dominant sans prétendre épuiser tous les usages d'un topic.

\subsection{Des topics aux familles thématiques}

Afin de rendre les résultats lisibles, les 17 topics ont été réunis a posteriori en six familles thématiques. Cette opération ne modifie ni le modèle ni l'affectation des segments : elle rapproche seulement des topics voisins pour organiser l'analyse et les visualisations temporelles. Il faut donc conserver deux niveaux de lecture : la famille fait apparaître un domaine général, tandis que les topics qui la composent en décrivent des dimensions parfois très différentes.

\begin{table}[H]
    \centering
    \caption{Regroupement interprétatif des 17 topics en six familles}
    \label{tab:familles-topics-bertopic}
    \begin{tabularx}{\textwidth}{
        >{\raggedright\arraybackslash}X
        >{\centering\arraybackslash}p{2.5cm}
        >{\centering\arraybackslash}p{2.1cm}
        >{\centering\arraybackslash}p{2.1cm}}
        \toprule
        \textbf{Famille thématique} & \textbf{Topics} & \textbf{Segments} & \textbf{Part des segments assignés} \\
        \midrule
        Corps, apparences et vie privée ou mondaine & 5, 8, 10 & \(15\,460\) & 25,17\,\% \\
        Travail, argent et conditions matérielles & 4, 6, 9, 14 & \(14\,814\) & 24,12\,\% \\
        Institutions, religion et pouvoir & 1, 12, 13, 15 & \(12\,107\) & 19,71\,\% \\
        Relations affectives, couple et famille & 2, 7, 11, 16 & \(10\,986\) & 17,89\,\% \\
        Espaces, paysages et architecture & 0 & \(5\,397\) & 8,79\,\% \\
        Guerre et violence collective & 3 & \(2\,655\) & 4,32\,\% \\
        \bottomrule
    \end{tabularx}
\end{table}

Ces proportions décrivent la taille des regroupements produits par le modèle, non l'importance intrinsèque des thèmes dans l'œuvre. Elles dépendent notamment de la segmentation, de la granularité choisie et de la capacité de certains lexiques à former un groupe homogène. Le tableau~\ref{tab:description-topics-bertopic} revient donc au niveau des topics et associe à chacun un intitulé analytique, son effectif et cinq de ses termes les plus discriminants.

{\small
\begin{longtable}{
    @{}
    >{\centering\arraybackslash}p{0.07\textwidth}
    >{\centering\arraybackslash}p{0.10\textwidth}
    p{0.33\textwidth}
    p{0.39\textwidth}
    @{}
}
\caption{Interprétation synthétique des 17 topics du modèle final}
\label{tab:description-topics-bertopic}\\
\toprule
\textbf{Topic} & \textbf{Effectif} & \textbf{Intitulé analytique} & \textbf{Termes discriminants}\\
\midrule
\endfirsthead
\multicolumn{4}{c}{\tablename\ \thetable{} --- suite}\\
\toprule
\textbf{Topic} & \textbf{Effectif} & \textbf{Intitulé analytique} & \textbf{Termes discriminants}\\
\midrule
\endhead
\midrule
\multicolumn{4}{r}{\textit{Suite à la page suivante}}\\
\endfoot
\bottomrule
\endlastfoot
\multicolumn{4}{@{}l}{\textbf{Corps, apparences et vie privée ou mondaine}}\\*
5  & \(5\,277\) & Nuit, sommeil et espace intime & oreiller, insomnie, recoucher, respiration, parquet\\
8  & \(7\,588\) & Vie domestique et sociabilité mondaine & thé, loge, cocher, piano, vestibule\\
10 & \(2\,595\) & Corps, vêtements et portraits physiques & hanche, blancheur, chevelure, corsage, physionomie\\
\addlinespace
\multicolumn{4}{@{}l}{\textbf{Travail, argent et conditions matérielles}}\\*
4  & \(3\,817\) & Finance, propriété et spéculation & banque, capital, spéculation, immeuble, faillite\\
6  & \(5\,714\) & Alimentation et sociabilité populaire & soupe, Gervaise, blanchisseur, zingueur, casserole\\
9  & \(2\,873\) & Machines, industrie et infrastructures & rail, berline, câble, fonte, acier\\
14 & \(2\,410\) & Dette et économie quotidienne & centime, rembourser, emprunter, usurier, créance\\
\addlinespace
\multicolumn{4}{@{}l}{\textbf{Institutions, religion et pouvoir}}\\*
1  & \(4\,554\) & Discours socio-religieux et projet collectif & nation, catholicisme, démocratie, salariat, solidarité\\
12 & \(2\,631\) & Hiérarchie ecclésiastique & éminence, prélat, monsignor, pape, congrégation\\
13 & \(1\,923\) & Enquête, preuve et justice & paraphe, accusé, arrestation, expert, audience\\
15 & \(2\,999\) & Conflit clérical et politique & jésuite, clérical, laïque, réactionnaire, municipal\\
\addlinespace
\multicolumn{4}{@{}l}{\textbf{Relations affectives, couple et famille}}\\*
2  & \(3\,188\) & Parole amoureuse et affective & pleure, mienne, mien, épouse, désire\\
7  & \(4\,725\) & Couple, maternité et sacrifice & épouse, maternité, virginité, union, sein\\
11 & \(1\,734\) & Supplication, pardon et devoir moral & supplie, conjure, miséricorde, confesse, pénitence\\
16 & \(1\,339\) & Union, séparation et intimité & séparation, union, fiançaille, fiancé, embrassement\\
\addlinespace
\multicolumn{4}{@{}l}{\textbf{Espaces, paysages et architecture}}\\*
0  & \(5\,397\) & Paysages naturels et espaces construits & toiture, feuillage, nef, coteau, dôme\\
\addlinespace
\multicolumn{4}{@{}l}{\textbf{Guerre et violence collective}}\\*
3  & \(2\,655\) & Guerre, armée et affrontement collectif & troupe, barricade, batterie, allemand, artillerie\\
\end{longtable}
}

\subsection{Lecture qualitative des six familles}

\paragraph{Corps, apparences et vie privée ou mondaine.} Cette première famille est la plus volumineuse, mais son unité tient moins au seul corps qu'aux scènes où celui-ci est observé dans un espace social ou intime. Le topic~5 rassemble la chambre, le sommeil, l'insomnie, la respiration et parfois la défaillance physique ; il est particulièrement présent dans \textit{Thérèse Raquin} et \textit{La Joie de vivre}. Le topic~10 correspond plus nettement au portrait : le corps y devient visible par la chevelure, les vêtements, les formes et la physionomie. Le topic~8, le plus grand du modèle avec \(7\,588\) segments, déplace l'attention vers les pratiques domestiques et mondaines, les visites, le théâtre, les moyens de transport et le personnel de maison. Sa forte représentation dans \textit{Nana}, \textit{Pot-Bouille} et \textit{Une page d'amour} montre que le modèle saisit moins une catégorie sociale unique qu'une mise en scène répétée des individus à travers leurs intérieurs, leurs apparences et leurs codes de sociabilité.

\paragraph{Travail, argent et conditions matérielles.} Les quatre topics de cette famille décrivent plusieurs échelles de la vie économique. Le topic~6 part du quotidien populaire : les repas, les métiers et les échanges familiers, fortement associés à \textit{L'Assommoir}, \textit{La Terre} et \textit{Le Ventre de Paris}. La présence de Gervaise parmi ses termes les plus discriminants montre toutefois qu'une partie de sa cohérence provient directement de \textit{L'Assommoir} et invite à ne pas l'assimiler sans réserve à l'ensemble des milieux populaires. Le topic~9 isole au contraire la matérialité industrielle par les rails, les câbles, les charpentes et les machines ; \textit{Germinal} en constitue le foyer principal. Le topic~4 porte sur la banque, la propriété et la spéculation, avec une concentration attendue dans \textit{L'Argent}, mais aussi dans \textit{La Curée} et \textit{Au Bonheur des dames}. Enfin, le topic~14 rend visible une économie plus diffuse, faite de centimes, d'emprunts, de dots et de créances. Leur regroupement met ainsi en relation le budget domestique, la consommation, le travail industriel et la finance : l'argent n'apparaît pas comme un thème isolé, mais comme un système de contraintes qui relie les corps, les métiers et les positions sociales.

\paragraph{Institutions, religion et pouvoir.} Cette famille associe quatre formes de discours institutionnel. Le topic~12 représente la hiérarchie catholique à travers ses titres, ses audiences et ses congrégations ; il est particulièrement caractéristique de \textit{Rome}, tout en étant déjà visible dans \textit{La Conquête de Plassans} et \textit{Son Excellence Eugène Rougon}. Le topic~15 porte davantage sur l'affrontement entre cléricaux, laïques, réactionnaires et socialistes, notamment dans \textit{Vérité}. Le topic~13 décrit le fonctionnement de la justice par l'enquête, les expertises, les témoignages et l'accusation ; sa concentration dans \textit{Vérité} confirme la cohérence de cette interprétation. Plus abstrait, le topic~1 rassemble les discours sur la nation, le catholicisme, la démocratie, le salariat et la transformation sociale. L'ensemble fait apparaître le passage des institutions comme cadres du récit aux institutions comme objets explicites de critique et de projet politique.

\paragraph{Relations affectives, couple et famille.} Les topics~2, 7, 11 et~16 décrivent l'intimité à partir de positions d'énonciation et de situations relationnelles différentes. Le topic~2 est dominé par l'adresse à l'autre, le désir et la plainte ; le topic~11 par la supplication, le pardon et l'obligation morale. Les topics~7 et~16 structurent plus directement le couple autour de la maternité, de la virginité, de l'union, de la séparation et des fiançailles. La famille est très présente dans les premiers romans, notamment \textit{La Confession de Claude}, \textit{Le Vœu d'une morte}, \textit{Thérèse Raquin} et \textit{Madeleine Férat}, mais elle ne relève pas d'un simple registre sentimental : les documents représentatifs associent régulièrement l'affect aux normes conjugales, au sacrifice, à la sexualité et au devoir familial.

\paragraph{Espaces, paysages et architecture.} Le topic~0 réunit des éléments naturels et construits : feuillages, coteaux et rives voisinent avec les toitures, les nefs, les dômes et les alignements. Cette proximité montre que le topic isole avant tout une opération descriptive fondée sur l'organisation visuelle de l'espace. Il est particulièrement présent dans \textit{La Faute de l'abbé Mouret}, \textit{Le Ventre de Paris} et \textit{Rome}, trois œuvres qui mobilisent pourtant des milieux très différents. Le regroupement est donc cohérent sur le plan discursif sans correspondre à un décor unique ; il peut contribuer à l'étude naturaliste du milieu, mais ne constitue pas à lui seul un marqueur exclusif du naturalisme.

\paragraph{Guerre et violence collective.} Le topic~3 est le plus spécialisé du modèle. Son vocabulaire militaire --- troupes, divisions, batteries, obus, artillerie et nationalités combattantes --- renvoie principalement à \textit{La Débâcle}, qui réunit \(1\,095\) de ses \(2\,655\) segments. Des concentrations plus modestes apparaissent dans \textit{La Fortune des Rougon}, \textit{Les Mystères de Marseille} et \textit{Germinal}, où le lexique de l'affrontement collectif peut être mobilisé dans des contextes insurrectionnels ou sociaux. La cohérence de ce topic facilite son interprétation, mais sa forte dépendance à quelques romans rappelle qu'un groupe lexical très net peut représenter un épisode spécialisé plutôt qu'une préoccupation uniformément distribuée dans toute l'œuvre.

\subsection{Évolution des topics au fil du temps}

Le modèle a été construit sans intégrer la date de publication. La fonction \path{topics_over_time} répartit ensuite les segments déjà assignés dans quinze intervalles chronologiques de largeur égale entre 1865 et 1903, puis recalcule une représentation lexicale pour chaque topic et chaque intervalle. Les figures ci-dessous ne montrent pas une courbe additionnée par famille : elles mettent en évidence, dans une même image, les trajectoires individuelles des topics qui composent cette famille.

L'option \texttt{normalize\_frequency=False} ayant été conservée dans le notebook, l'axe vertical représente des effectifs bruts. Or le nombre de segments assignés varie de \(1\,340\) à \(6\,307\) selon les intervalles, soit un rapport de 1 à 4,7. Une courbe élevée peut donc traduire à la fois la place d'un topic dans les romans concernés et un volume textuel plus important. Pour limiter ce biais, l'analyse confronte les effectifs visibles à la part qu'ils représentent dans chaque intervalle et à leur répartition par roman ; elle porte sur la fréquence de topics globaux, non sur une évolution autonome de leur sens.

Les abscisses correspondent aux bornes inférieures de classes d'environ 2,53 années, et non à des dates d'observation exactes ; le point situé vers 1890,3 regroupe par exemple les romans publiés en 1891 et 1892. Les lignes reliant les points facilitent la lecture, mais ne doivent pas être comprises comme une continuité annuelle entre des œuvres publiées de manière discontinue.

\subsubsection{Corps, apparences et vie privée ou mondaine}

\begin{figure}[H]
    \centering
    \includegraphics[width=\textwidth,trim=0 0 0 75,clip]{\detokenize{images/chapitre_2/visualisation_temps/Corps, apparences et vie privée ou mondaine.png}}
    \caption{Évolution en effectifs bruts des topics 5, 8 et 10 : corps, apparences et vie privée ou mondaine}
    \label{fig:temps-corps-vie-privee}
\end{figure}

La figure~\ref{fig:temps-corps-vie-privee} fait surtout apparaître le poids du topic~8. Celui-ci atteint 978 segments dans l'intervalle contenant les publications de 1878 et 1880, puis 586 dans celui de 1882 ; ces valeurs sont principalement portées par \textit{Une page d'amour}, \textit{Nana} et \textit{Pot-Bouille}. Les trois topics réunis représentent respectivement 46,6\,\% et 45,7\,\% des segments assignés dans ces deux intervalles, ce qui confirme que les sommets ne résultent pas uniquement du volume du corpus. Le pic brut le plus élevé de la famille apparaît toutefois en 1883--1885 avec \(1\,737\) segments : ramené aux \(6\,307\) segments de l'intervalle, il ne représente plus que 27,5\,\%. La distinction entre effectif et proportion modifie donc sensiblement la lecture de la courbe.

Les topics~5 et~10 suivent des trajectoires moins concentrées. Le premier atteint 696 segments en 1883--1885 et relie des scènes nocturnes ou intimes présentes dans plusieurs univers romanesques ; le second reste plus faible et diffuse les descriptions corporelles sur l'ensemble de la chronologie. Cette famille ne dessine donc pas une progression continue. Elle met plutôt en évidence un moment de forte densité domestique et mondaine autour de 1878--1882, entouré d'une présence plus régulière du corps, du sommeil et du portrait.

\subsubsection{Travail, argent et conditions matérielles}

\begin{figure}[H]
    \centering
    \includegraphics[width=\textwidth,trim=0 0 0 75,clip]{\detokenize{images/chapitre_2/visualisation_temps/Travail, argent et conditions matérielles.png}}
    \caption{Évolution en effectifs bruts des topics 4, 6, 9 et 14 : travail, argent et conditions matérielles}
    \label{fig:temps-travail-argent}
\end{figure}

Les courbes de la figure~\ref{fig:temps-travail-argent} montrent que la famille matérielle ne possède pas un sommet unique. Le topic~6 atteint 1\,253 segments dans l'intervalle 1876--1877, dont \(1\,138\) dans \textit{L'Assommoir} ; avec les topics~4, 9 et~14, la famille représente alors 44,6\,\% des segments assignés. En 1883--1885, le topic~9 culmine à 684 segments, dont 496 dans \textit{Germinal}, tandis que le topic~6 demeure élevé sous l'effet conjoint d'\textit{Au Bonheur des dames}, de \textit{La Joie de vivre} et de \textit{Germinal}. L'intervalle rassemble \(2\,348\) segments de la famille, soit 37,2\,\% des segments qui y sont assignés.

Le topic~4 connaît son maximum en 1891--1892 avec 594 segments, dont 498 appartiennent à \textit{L'Argent}. Le topic~14 ne présente pas une spécialisation comparable : son vocabulaire de la dette et des petites sommes demeure réparti entre de nombreux romans. La succession de ces pics n'indique donc pas le remplacement d'un thème par un autre. Elle montre plutôt trois formes de détermination matérielle que le modèle distingue : la vie populaire et alimentaire autour de \textit{L'Assommoir}, la mécanisation du travail autour de \textit{Germinal}, puis la finance autour de \textit{L'Argent}. C'est cette articulation, plus que la hauteur globale d'une courbe, qui rend la famille centrale pour l'analyse du projet naturaliste.

\subsubsection{Institutions, religion et pouvoir}

\begin{figure}[H]
    \centering
    \includegraphics[width=\textwidth,trim=0 0 0 75,clip]{\detokenize{images/chapitre_2/visualisation_temps/Institutions, religion et pouvoir.png}}
    \caption{Évolution en effectifs bruts des topics 1, 12, 13 et 15 : institutions, religion et pouvoir}
    \label{fig:temps-institutions-pouvoir}
\end{figure}

La figure~\ref{fig:temps-institutions-pouvoir} présente l'inflexion chronologique la plus nette. À l'échelle des ensembles romanesques, la famille représente 12,2\,\% des segments des \textit{Rougon-Macquart}, contre 37,8\,\% dans les \textit{Trois Villes} et 37,0\,\% dans les trois romans publiés du cycle des \textit{Quatre Évangiles} présents dans le corpus. Dans l'intervalle qui contient uniquement \textit{Rome} en 1896, les quatre topics réunissent \(1\,674\) segments, soit 47,7\,\% du roman. Le dernier intervalle, qui regroupe \textit{Travail} et \textit{Vérité}, atteint \(2\,831\) segments et 45,4\,\% des segments assignés.

Cette hausse ne repose pas sur un seul topic. Le topic~1, consacré aux discours socio-religieux et aux projets collectifs, atteint 1\,349 segments dans le dernier intervalle ; le topic~15, centré sur le conflit clérical et politique, en compte 676 ; le topic~13, judiciaire, 568. Le topic~12 suit une trajectoire plus discontinue : il est déjà élevé dans l'intervalle 1873--1875, notamment à cause de \textit{La Conquête de Plassans}, puis atteint 450 segments dans \textit{Rome}. Les résultats ne montrent donc pas une apparition tardive des institutions, mais une reconfiguration : les œuvres finales combinent plus fortement discours idéologique, conflit religieux, organisation politique et procédure judiciaire.

\subsubsection{Relations affectives, couple et famille}

\begin{figure}[H]
    \centering
    \includegraphics[width=\textwidth,trim=0 0 0 75,clip]{\detokenize{images/chapitre_2/visualisation_temps/Relations affectives, couple et famille.png}}
    \caption{Évolution en effectifs bruts des topics 2, 7, 11 et 16 : relations affectives, couple et famille}
    \label{fig:temps-relations-famille}
\end{figure}

La figure~\ref{fig:temps-relations-famille} révèle une présence particulièrement forte au début du corpus. La famille représente 33,5\,\% des segments de l'intervalle 1865--1867 et 51,1\,\% de celui qui contient \textit{Madeleine Férat} en 1868. À l'échelle des romans, elle atteint également 59,2\,\% dans \textit{La Confession de Claude}, 44,5\,\% dans \textit{Le Vœu d'une morte} et 38,9\,\% dans \textit{Thérèse Raquin}. Ces valeurs confirment que les premiers récits structurent fortement les affects autour de la parole amoureuse, de l'union, de la jalousie, du sacrifice et de la faute.

La famille ne disparaît cependant pas avec le développement des grands romans sociaux. Sa part tombe à 6,2\,\% en 1876--1877, puis fluctue avant d'atteindre 21,8\,\% en 1898--1899 et 19,4\,\% en 1901--1903. Le topic~7, centré sur le couple et la maternité, porte une grande partie de cette reprise, notamment dans \textit{Fécondité}. La courbe décrit donc moins un recul linéaire du registre affectif qu'un changement de contexte : d'abord au premier plan de récits intimes, il est ensuite réinscrit dans des ensembles où la famille est articulée à des enjeux sociaux, moraux ou collectifs.

\subsubsection{Espaces, paysages et architecture}

\begin{figure}[H]
    \centering
    \includegraphics[width=\textwidth,trim=0 0 0 75,clip]{\detokenize{images/chapitre_2/visualisation_temps/Espaces, paysages et architecture.png}}
    \caption{Évolution en effectifs bruts du topic 0 : espaces, paysages et architecture}
    \label{fig:temps-espaces-paysages}
\end{figure}

Le topic~0 est présent dans les quinze intervalles, mais sa courbe reste irrégulière. Son maximum brut, 742 segments en 1873--1875, est exactement formé par \textit{Le Ventre de Paris}, \textit{La Conquête de Plassans} et \textit{La Faute de l'abbé Mouret}. Ces œuvres associent respectivement architecture marchande, ville provinciale et paysage végétal, ce qui confirme que le topic repère une manière de construire l'espace plutôt qu'un lieu homogène. Un second sommet de 639 segments correspond à \textit{Rome} en 1896 ; il représente alors 18,2\,\% du roman, soit la part la plus élevée de la chronologie.

La présence transversale du topic montre que la description du milieu accompagne plusieurs phases de l'œuvre, mais ses variations sont largement dépendantes des décors propres à chaque roman. On ne peut donc pas déduire de la courbe une intensification ou un effacement général de la description. Elle indique plutôt que BERTopic rapproche des passages où l'environnement est organisé comme un ensemble visible, qu'il s'agisse d'un paysage, d'un édifice ou d'un espace urbain.

\subsubsection{Guerre et violence collective}

\begin{figure}[H]
    \centering
    \includegraphics[width=\textwidth,trim=0 0 0 75,clip]{\detokenize{images/chapitre_2/visualisation_temps/Guerre et violence collective.png}}
    \caption{Évolution en effectifs bruts du topic 3 : guerre et violence collective}
    \label{fig:temps-guerre-violence}
\end{figure}

La figure~\ref{fig:temps-guerre-violence} se distingue par un pic massif et ponctuel. Le topic~3 atteint \(1\,143\) segments dans l'intervalle 1891--1892, soit 23,1\,\% des segments assignés de cette période. Parmi eux, \(1\,095\), ou 95,8\,\%, proviennent de \textit{La Débâcle} ; ce roman rassemble à lui seul 41,2\,\% de tous les segments du topic. L'association entre le vocabulaire militaire et cette œuvre est donc beaucoup plus directe que pour les familles transversales.

Les autres hausses restent nettement plus faibles : 234 segments dans l'intervalle de 1871, principalement sous l'effet de \textit{La Fortune des Rougon}, et 245 en 1883--1885, dont 193 dans \textit{Germinal}. Ces occurrences montrent que l'affrontement collectif traverse aussi les scènes insurrectionnelles ou sociales, mais elles ne remettent pas en cause la spécialisation du topic. Cette trajectoire constitue un cas utile : elle illustre la capacité du modèle à isoler un lexique historiquement et narrativement concentré, tout en rappelant qu'un sommet temporel peut d'abord être l'empreinte d'un seul roman.

\subsection{Bilan interprétatif}

Les six familles ne suivent pas une évolution commune. Les relations affectives dominent plusieurs œuvres de jeunesse puis sont reformulées dans les romans tardifs ; le travail et les conditions matérielles connaissent des concentrations successives autour de la vie populaire, de l'industrie et de la finance ; les institutions deviennent beaucoup plus présentes dans les cycles finaux ; enfin, la guerre dépend très fortement de \textit{La Débâcle}. Les espaces et les scènes privées restent plus transversaux, mais leur fréquence varie avec les milieux particuliers mis en récit.

Cette lecture confirme l'intérêt de BERTopic pour identifier des régularités de contenu à l'échelle d'un vaste corpus. Le modèle distingue notamment les mécanismes économiques, les univers professionnels, les institutions, les formes de sociabilité et l'inscription des personnages dans des espaces matériels. Ces résultats rejoignent plusieurs objets privilégiés par l'enquête naturaliste, sans suffire à définir le naturalisme dans son ensemble : aucun topic n'isole clairement l'hérédité, et les courbes ne mesurent ni la construction narrative, ni la causalité entre milieu et personnage, ni le traitement stylistique des scènes.

Le regroupement en familles doit enfin rester réversible. Il facilite la comparaison temporelle, mais il peut masquer des trajectoires opposées à l'intérieur d'un même ensemble, par exemple entre l'économie domestique du topic~14 et la finance du topic~4, ou entre la hiérarchie ecclésiastique du topic~12 et la justice du topic~13. L'apport principal de l'analyse réside donc moins dans une taxonomie définitive que dans l'articulation de trois échelles : les segments représentatifs donnent accès au contexte, les topics isolent des régularités sémantiques et les familles permettent de suivre leurs recompositions dans l'œuvre.
